% !TeX root = ../main.tex
\section{Постановка задачи и исходные данные}
Рассматривается задача слежения за доплеровской частотой $\Omega(t) = \frac{\omega_0}{c} V(t)$ принимаемого радиосигнала на несущей частоте $f_0 = 1602\text{ МГц}$ ($\omega_0 = 2\pi f_0$, $c = 3 \cdot 10^8\text{ м/с}$). Динамика изменения доплеровской частоты и скорости ее изменения $v(t) = \frac{\omega_0}{c} a(t)$ описывается стохастическими дифференциальными уравнениями
\begin{equation}
\frac{d\Omega}{dt} = v(t), \qquad \frac{dv}{dt} = -\alpha v(t) + \alpha \xi(t),
\end{equation}
где $\alpha = 1\text{ с}^{-1}$ --- ширина спектра флуктуаций радиального ускорения, $\xi(t)$ --- формирующий белый гауссовский шум (БГШ) со спектральной плотностью мощности (СПМ)
\begin{equation}
S_\xi = 2 \sigma_a^2 \alpha \left(\frac{\omega_0}{c}\right)^2,
\end{equation}
где $\sigma_a$ --- среднеквадратическое значение радиального ускорения.

Эквивалентные наблюдения на выходе линеаризованного частотного дискриминатора записываются в виде
\begin{equation}
\tilde{y}(t) = \Omega(t) + \tilde{n}(t),
\end{equation}
где $\tilde{n}(t)$ --- эквивалентный БГШ дискриминатора с нулевым математическим ожиданием и односторонней СПМ
\begin{equation}
\tilde{N}_0(q_{c/n_0}) = \frac{2}{q_{c/n_0} T^2} \left(1 + \frac{1}{2 q_{c/n_0} T}\right),
\end{equation}
при интервале накопления $T = 10\text{ мс}$ и отношении сигнал/шум $q_{c/n_0} = 10^{0.1 \cdot q_{c/n_0}[\text{дБГц}]}$.

В дискретном времени с шагом $T = 0.01\text{ с}$ модель пространства состояний и наблюдений представляется системой разностных уравнений
\begin{align}
\Omega_k &= \Omega_{k-1} + v_{k-1} T, \\
v_k &= v_{k-1}(1 - \alpha T) + \alpha T \xi_{k-1}, \\
\tilde{y}_k &= \Omega_k + \tilde{n}_k,
\end{align}
где $\xi_{k-1}$ и $\tilde{n}_k$ --- независимые дискретные БГШ с дисперсиями $\sigma_\xi^2 = \frac{S_\xi}{2T}$ и $\sigma_n^2 = \frac{\tilde{N}_0(q_{c/n_0})}{2T}$ соответственно.
